%======================================================================
\chapter{Introduction}\label{ch:intro}
%======================================================================

\section{What this volume is about}

Volume~4 developed the plane topology that can be done with logarithms and polygons. This volume
adds the three remaining pieces used by the classification of homogeneous plane continua: the
Jordan curve theorem for arbitrary simple closed curves, R.~L.~Moore's axiomatic characterization of
the plane, and Moore's theorem that the space of pieces of a suitable decomposition of the plane into
continua is again a plane. The last is the form in which the volume is used later: it turns a
monotone image of a plane continuum into a plane set.

The volume follows two papers of Moore, and the order of the formal proof:
\emph{Concerning upper semi-continuous collections of continua} (1925), and \emph{Concerning a set of
postulates for plane analysis situs} (1919), which rests on his \emph{On the foundations of plane
analysis situs} (1916, cited as F.A.). The 1925 paper proves its main theorem by verifying the
axioms of the 1919 paper; on the way, its Theorem~11 is a Jordan curve theorem for curves of
elements, which for the decomposition into points is the classical one.

The volume has nine chapters.

\begin{itemize}
\item \textbf{Chapter~\ref{ch:arcs}: Arcs and simple closed curves.} Irreducible continua between two
points with no other non-cut points are arcs; the cut-point characterizations of arcs and circles;
limits of arcs; replacing a piece of an arc.
\item \textbf{Chapter~\ref{ch:usc}: Decompositions of the plane into non-separating continua.} Moore
decompositions, the quotient space of elements, its metrizability, the correspondence of components,
and Moore's Theorems~1--4.
\item \textbf{Chapter~\ref{ch:elements}: Arcs of elements.} Moore's Theorems~5--10 and~12, and the
decomposition into points.
\item \textbf{Chapter~\ref{ch:jordan}: Mullikin's theorem and the Jordan curve theorem.} Mullikin's
theorem on two non-separating continua with disconnected intersection, Moore's Theorem~11, and the
Jordan curve theorem.
\item \textbf{Chapter~\ref{ch:axioms}: Moore's axioms and the bridge to topology.} The eight axioms,
the topology generated by the regions, metrizability, arcs in domains (F.A.\ Theorem~15).
\item \textbf{Chapter~\ref{ch:curves}: Closed curves in a Moore space.} F.A.\ Theorems~20--43:
interiors and exteriors, cross-cuts, closed curves in regions, chords.
\item \textbf{Chapter~\ref{ch:rulings}: Double rulings and Moore's Theorem~D.} Quadrilaterals, double
rulings, a corrected Theorem~C, and fine rulings of a quadrilateral.
\item \textbf{Chapter~\ref{ch:plane}: Grids, the Schoenflies theorem and Moore's characterization of the
plane.} A compact set with a grid is a square; the Schoenflies theorem and the annulus lemma in a
Moore space; exhaustion by closed curves; every Moore space is a plane.
\item \textbf{Chapter~\ref{ch:quotient}: The space of elements is a plane.} Moore's Theorems~13--22,
monotone decompositions of a plane continuum, and the cross-cut, Schoenflies and collar theorems in
the plane.
\end{itemize}

The main results are:

\begin{theorem}[Jordan curve theorem; Theorem~\ref{thm:jct}]
If $J\subseteq\R^2$ is homeomorphic to a circle, $\R^2\setminus J$ has exactly two components, one bounded
and one unbounded, and $J$ is the frontier of each.
\end{theorem}

\begin{theorem}[Moore's characterization of the plane; Theorem~\ref{thm:moore-plane}]
A set with a class of regions satisfying Moore's Axioms~1--8 is homeomorphic to $\R^2$, with the
topology generated by the regions.
\end{theorem}

\begin{theorem}[Schoenflies; Theorems~\ref{thm:schoenflies} and~\ref{thm:plane-schoenflies}]
The closure of the bounded complementary domain of a simple closed curve in the plane is a closed disc
(a square), by a homeomorphism that can be prescribed on one arc of the curve.
\end{theorem}

\begin{theorem}[Moore 1925; Theorem~\ref{thm:m25-22}]
If the plane is decomposed into compact connected sets, none of which separates the plane, in an upper
semicontinuous way, then the space of these sets with the quotient topology is homeomorphic to $\R^2$.
\end{theorem}

\begin{corollary}[Corollary~\ref{cor:monotone-plane}]
A monotone image of a plane continuum whose fibres do not separate the plane is homeomorphic to a compact
subset of the plane.
\end{corollary}

\section{Where these results are used}

\begin{itemize}
\item Moore's Theorem~22 and the monotone-image corollary are used in Hagopian's arguments on
indecomposable homogeneous plane continua (Volume~6), where a plane continuum is decomposed into layers or
into the fibres of a monotone map, and in Jones's work on decomposable homogeneous continua (Volume~7):
in each case the space of pieces is recognised as a plane, or as a plane continuum.
\item The Jordan curve theorem and the circle characterization (Chapter~\ref{ch:arcs}) are used throughout
Hagopian's constructions and in Jones's study of circle-like continua (Volumes~6 and~7).
\item The cross-cut, Schoenflies and collar theorems in the plane (Chapter~\ref{ch:quotient}) are used in
Hagopian's constructions (Volume~6) and in Jones's arguments (Volume~7).
\end{itemize}

\section{Where the formal proof differs from the literature}\label{sec:intro-diff}

\begin{enumerate}
\item \textbf{Topology instead of Moore's primitive notions.} Moore works with limit elements, domains and
regions defined by his axioms. The formal proof translates them into the quotient topology
(Proposition~\ref{prop:moore-notions}) and into the topology generated by the regions
(Proposition~\ref{prop:bridge}), and then uses standard topology: metrizability, compactness, the
cut-point characterizations of arcs and circles.
\item \textbf{Printed statements that are false.} Several statements are false as printed and are proved here in
a corrected form; in each case the formal development proves the negation of the printed form.
\begin{itemize}
\item Moore 1925, Theorem~2, with ``complementary domain'' in Moore's sense (Remark~\ref{rem:m25-2-printed}).
\item Moore 1925, Theorems~10 and~11, with Moore's printed definition of a simple closed curve of
elements, and Theorem~7, with ``$p$'' for ``$x$'' (Chapter~\ref{ch:elements}, after
Definition~\ref{def:arc-elt} and after Theorem~\ref{thm:m25-5-10}).
\item Mullikin's theorem as quoted by Moore (Remark~\ref{rem:mullikin-printed}).
\item F.A.\ Theorems~4, 14, 19 and~35, which fail for the empty region, and Theorem~43 for an empty family
(Remark~\ref{rem:empty-region}, Theorems~\ref{thm:fa35-36} and~\ref{thm:fa43}).
\item Moore 1919, Theorem~C (Theorem~\ref{thm:theoremC}): a curve may lie within the quadrilateral without
crossing it suitably. The formal proof of Theorem~D uses a corrected Theorem~C for curves that cross the
quadrilateral transversally.
\end{itemize}
\item \textbf{The Jordan curve theorem without the characterization of the plane.} The Jordan curve theorem is
obtained from Moore's Theorem~11 for the decomposition into points (Chapter~\ref{ch:jordan}), using
Janiszewski's and Borsuk's theorems from Volume~4 but none of Chapters~\ref{ch:axioms}--\ref{ch:plane}.
\item \textbf{Coordinates by a grid.} Moore 1919 reaches the plane through ``nets'' of curves (Theorems~E--G).
The formal proof proves a general recognition theorem for squares (a compact set with two families of grid
lines, Theorem~\ref{thm:grid-square}), derives the Schoenflies theorem in a Moore space from Theorem~D, and
pastes squares and collars along an exhaustion (Chapter~\ref{ch:plane}); Theorems~E and~G become corollaries.
\item \textbf{Moore's Axioms~6 and~7 for the space of elements.} Moore obtains them through F.A.\ and the 1919
paper; here they are proved directly from Theorems~12, 15, 18 and~19 (Theorem~\ref{thm:m25-21}).
\item \textbf{Proofs in outline.} Chapters~\ref{ch:curves}--\ref{ch:plane} contain some long case analyses
(several of F.A.'s theorems on closed curves, the corrected Theorem~C, Theorem~D). For these the notes give
the proofs in outline, marked as such, and indicate which earlier results they use; the complete arguments are
in the formal development, at the declarations named in the ``Formal counterpart'' lines.
\end{enumerate}

\section{Prerequisites and conventions}\label{sec:intro-conv}

The reader should know the basic metric topology of \vone{Chapter~1} and the results of Volume~4 on
logarithms, Brouwer's theorem, Janiszewski's theorem and Borsuk's theorem. From Volume~3 we use
irreducible continua, non-cut points and upper semicontinuous decompositions.

\begin{itemize}
\item \textbf{The plane.} As in Volume~4, $\R^2$ carries the sup metric, so that balls are squares.
\item \textbf{Elements.} For a decomposition of the plane, $\el x$ is the element containing $x$, $G$ is the space of
elements, and $\pts H$ is the point set of a set $H$ of elements (Definition~\ref{def:moore-decomp}).
\item \textbf{Arcs and curves.} An \emph{arc} from $a$ to $b$ is the image of an embedding of $[0,1]$ with $0\mapsto a$,
$1\mapsto b$; a \emph{simple closed curve} is the union of two arcs with only their ends in common, equivalently
a set homeomorphic to a circle (Chapter~\ref{ch:arcs}). $\oarc{\gamma}$ is the arc minus its ends.
\item \textbf{Moore spaces.} In Chapters~\ref{ch:axioms}--\ref{ch:plane}, $S$ is a set with regions satisfying some of
Axioms~1--8; $I(J)$ is the interior of a closed curve $J$, and ``within'', ``without'' mean in the interior,
in the exterior.
\item \textbf{Frontier.} $\fr A$ is the topological frontier; $\bd M$ is Moore's boundary.
\end{itemize}

\paragraph{Formal counterparts.} As in the earlier volumes, many results are followed by a line
``Formal counterpart:'' naming the corresponding declaration of the Lean development (all names
begin with \lean{HPC.}). These lines may be skipped. Appendix~\ref{app:formal} collects them.

\paragraph{Exercises.} Each chapter ends with exercises. None of them is used later.
